\noindent\textbf{Elección del procedimiento.} La semejanza exige igual Arquímedes y, para la misma rama de la tabla, igual Reynolds crítico. Con el mismo líquido, $U_{c2}d_2=U_{c1}d_1$. Primero identificar el ensayo cuyo $Bd^3$ coincide, y después escalar velocidad inversamente con diámetro.

\begin{enumerate}

\item Comparar este grupo con el del ensayo de 80 N/m$^3$ y diámetro 4 mm.

\[\fitmath{Ar2=\frac{\rho\cdot B2\cdot {d2}^{3}}{{\mu}^{2}}}\]

\[\fitmath{Ar2=\frac{(1000)\cdot (10)\cdot {(0.008)}^{3}}{{(0.001)}^{2}}\approx 5120\,}\]

\item Conservar el producto velocidad por diámetro.

\[\fitmath{Uc2=\frac{U1\cdot d1}{d2}}\]

\[\fitmath{Uc2=\frac{(0.0648)\cdot (0.004)}{(0.008)}\approx 0.0324\,\mathrm{m/s}}\]

\end{enumerate}
